\section{The Invariant: Machine-Checkable Security Claims}

\label{sec:invariant}

\begin{invariant}[Machine-checkable security]
Every security property of the Lux Network either:
\begin{enumerate}[leftmargin=2em]
  \item Has a machine-checkable proof (formal verification in Coq, Lean, or TLA+), or
  \item Has an automated test suite that verifies the property against known test
        vectors, or
  \item Is explicitly listed as an unverified assumption with a justification for
        why it is believed to hold.
\end{enumerate}
There is no fourth category. ``We believe it is secure'' without one of the above
annotations is not a valid security claim.
\end{invariant}

This invariant is not aspirational---it is a \emph{process requirement}. Every
protocol change must be accompanied by an update to one of the three categories.
New cryptographic primitives require either a proof reference or NIST KAT validation.
New consensus rules require either a formal proof or a TLA+ model. New implementations
require either formal verification or an automated test suite with coverage metrics.

The purpose of this invariant is not to achieve 100\% formal verification (which is
impractical for a large system). The purpose is to eliminate the category of
``security claims that nobody has checked.'' Every claim is checked by a machine, a
proof system, or a test suite---or it is explicitly marked as unchecked.

%% ========================================================================
%%  10. RELATED WORK
%% ========================================================================
